\section{Related work}
\label{2}
\subsection{Continual Learning}
\label{2.1}
Continual Learning (CL) aims to train models sequentially on evolving tasks while retaining prior knowledge. In the CL setting, new classes arrive without access to old data, often leading to catastrophic forgetting. To address this, several strategies have emerged.  

Kirkpatrick \textit{et al.} ~\cite{kirkpatrick2017overcoming} introduced Elastic Weight Consolidation (EWC), an early exemplar-free approach that adds a Fisher-based quadratic penalty to protect parameters critical for past tasks. Buzzega \textit{et al.} ~\cite{buzzega2020dark} proposed Dark Experience Replay++ (DER++), combining rehearsal with distillation by storing both samples and logits to stabilize decision boundaries. Caccia \textit{et al.} ~\cite{caccia2021new} developed ER-ACE and ER-AML where cross-entropy is applied asymmetrically, isolating new-class updates while replay consolidates all classes. Douillard \textit{et al.} ~\cite{douillard2022dytox} presented DyTox, a transformer-based method using task-specific tokens and a replay buffer to scale across tasks without task IDs. Cotogni \textit{ et al.} ~\cite{cotogni2025exemplar} introduced GCAB, which employs gated class attention and feature drift compensation for exemplar-free ViT learning. Finally, Mohamed \textit{et al.} ~\cite{mohamed2023d3former} proposed D3Former, which debiases logits and preserves attention maps to balance performance across old and new classes.
  
 
\subsection{Machine Unlearning}
\label{2.2}
Machine Unlearning (MU) has seen significant advancements, focusing on methods that effectively balance the removal of data influence and the preservation of retained knowledge. Kurmanji \textit{et al.} \cite{kurmanji2023towards} tackled the problem of unlearning by pushing the forget distribution toward a uniform distribution, while ensuring that the retain distribution follows the normal loss function. This approach ensures that models forget the target data without significantly compromising performance on the remaining data. Zhao \textit{ et al.} \cite{zhao2024continual} addressed the issue of continual forgetting in the context of continual learning by utilizing GS-LoRA, which targets the FFN modules with a group sparsity regularizer. Cha \textit{et al.} \cite{cha2024learning} introduced adversarial examples as a technique to ensure a cleaner and more robust removal of forget data. By generating adversarial examples, the model is encouraged to misclassify the forget data, making it harder for the model to remember or retain the forgotten samples. Fan \textit{et al.} \cite{fan2023salun} proposed a novel method that computes the weight saliency matrix for data to be forgotten. By identifying the most relevant weights for the forgotten samples, the model updates only those weights, which helps to avoid catastrophic forgetting and ensures that the model retains critical information from the retained data. Tarun \textit{et al.}~\cite{tarun2023fast} proposed a method that corrupts the knowledge to be forgotten by using noise and then performs repair work to maintain the retained knowledge. Panda \textit{et al.}~\cite{panda2024fast} proposed a weak unlearning approach for black-box GANs that filters undesired outputs by computing projection similarity between sampled latent vectors and a learned representation of unwanted features in the latent space. Wang \textit{et al.}~\cite{wang2025malicious} proposed MCC-Fed, which detects malicious clients via Euclidean distance-based deviation analysis and employs a Lipschitz-inspired contribution-aware metric as a regularization term to precisely unlearn their negative influence without requiring auxiliary datasets.

\subsection{Parameter-Efficient Fine-Tuning}
\label{2.3}
Fine-tuning large pretrained models such as vision and language models on downstream tasks has become a dominant paradigm in modern deep learning. Parameter-efficient fine-tuning techniques are widely adopted, as they substantially reduce trainable parameters without significant performance degradation. Addition-based approaches \cite{liu2024gpt} introduce new trainable components while freezing the original model; other works include \cite{houlsby2019parameter,li2021prefix}. Freezing-based techniques \cite{lee2019would} represent another approach, selectively updating only specific parameters or layers while keeping the rest frozen. Parameter factorization methods \cite{chavan2023one,valipour2022dylora} form the third category, decomposing weight updates into low-rank matrices to achieve efficiency. Among these, Low-Rank Adaptation (LoRA)~\cite{hu2022lora} has emerged as particularly effective, decomposing weight updates into low-rank matrices that can be efficiently trained and merged. Our work builds upon LoRA's foundation to address the unique challenges of continual learning-unlearning.

\subsection{Continual Learning-Unlearning (CLU)}
\label{2.4}
The existing CLU works, such as Shibata \textit{et al.} ~\cite{shibata2021learning}, Liu \textit{et al.}  ~\cite{liu2022continual}, Chatterjee \textit{et al.} ~\cite{chatterjee2024unified}, and Huang \textit{et al.} ~\cite{huang2025unified}, integrate CL and MU, providing theoretical foundations for adaptive knowledge management. While this field has gained attention, existing approaches still face significant design and efficiency challenges.  

Shibata \textit{et al.} ~\cite{shibata2021learning} first explored a CLU-like system using mnemonic codes, enabling selective forgetting by discarding class-specific codes and new learning by embedding fresh ones. However, this method requires retraining from scratch with mnemonic embeddings, is incompatible with pre-trained models, and doubles computational cost. Liu \textit{et al.} ~\cite{liu2022continual} proposed CLPU-DER++, which achieves exact unlearning via isolated temporary networks that can be deleted while retaining a permanent model. It guarantees privacy and supports selective forgetting without original data. Yet, it demands retraining from scratch, incurs exponential storage overhead with multiple tasks, and cannot revise permanent knowledge. Chatterjee \textit{et al.} ~\cite{chatterjee2024unified} introduced UniCLUN, a dual-teacher distillation framework with one CL teacher, one UL teacher, and a student model. It is the first unified CL–UL approach, handling mixed learning–forgetting sequences adaptively. Nevertheless, it incurs parameter overhead, requires heavy computation, and depends on replay buffers that raise privacy risks. Adhikari \textit{et al.} ~\cite{adhikari2025unlearning} proposed UnCLe, a hypernetwork-based framework that generates task-specific parameters conditioned on task embeddings, enabling data-free unlearning by aligning forget-task parameters with Gaussian noise. However, it supports only task-level unlearning and cannot be directly applied to pretrained models like DeiT, as it requires training a hypernetwork to generate entire model parameters from scratch rather than fine-tuning existing weights. More recently, Huang \textit{et al.} ~\cite{huang2025unified} proposed a gradient-based, task-agnostic CLU framework optimized via KL divergence, aiming for efficiency and adaptability. However, it requires computing online Hessian approximations through costly inner-loop optimization, lacks support for pretrained vision transformers, and has only been validated on models trained from scratch rather than leveraging existing foundation models.

Despite these contributions, none of the above methods are resource-friendly: most update nearly all parameters, making them inefficient for large-scale pre-trained models. Moreover, they have only been tested on limited scenarios involving a few additions and deletions, raising questions about their claims of `continual' operation. In contrast, our BID-LoRA is explicitly designed to be parameter-efficient, resource-conscious, and validated across long sequences of CLU requests, demonstrating genuine continual adaptability at scale.